% from ega4-2-fr.tex
\oldpage[IV-2]{153}
\item[c)] $\mathcal O_x$ est une $\mathcal O_y$-algèbre formellement lisse
(0, 19.3.1) pour les topologies préadiques sur $\mathcal O_x$ et $\mathcal O_y$.
\item[c')] $\mathcal O_x$ est une $\mathcal O_y$-algèbre formellement lisse
(0, 19.3.1) pour les topologies discrètes sur $\mathcal O_x$ et $\mathcal O_y$.
\end{enumerate}
\end{theorem}

\begin{proof}
L'équivalence de a) et b) résulte des définitions (6.8.1) et du fait que,
pour les préschémas localement noethériens, les morphismes localement de type
fini sont localement de présentation finie (1.4.2). En second lieu, pour que
$\mathcal O_x$ soit une $\mathcal O_y$-algèbre formellement lisse pour les
topologies préadiques, il faut et il suffit que $\mathcal O_x$ soit un
$\mathcal O_y$-module plat et que $\mathcal O_x\otimes_{\mathcal O_y}k(y)$
soit une $k(y)$-algèbre formellement lisse pour sa topologie préadique
(0, 19.7.1); mais pour que $\mathcal O_x\otimes_{\mathcal O_y}k(y)$ soit une $k(y)$-algèbre formellement lisse pour sa topologie préadique,
il faut et il suffit qu'elle soit une $k(y)$-algèbre géométriquement régulière
(0, 19.6.6); ceci démontre donc l'équivalence de b) et c). Enfin, pour prouver
l'équivalence de c) et c'), on peut se borner au cas où $Y=\operatorname{Spec}(A)$
et $X=\operatorname{Spec}(C)$ sont affines, $A$ étant noethérien et $C$ une
$A$-algèbre de type fini, que l'on peut donc écrire sous la forme
$A[T_1,\ldots,T_n]/\mathfrak J$. Comme ici $\mathfrak J/\mathfrak J^2$ est un
$C$-module de présentation finie, l'équivalence de c) et c') résulte de
(0, 22.6.4) appliqué à $A$, $B=A[T_1,\ldots,T_n]$ et $C=B/\mathfrak J$.
\end{proof}

\begin{corollary}[6.8.7]
\label{IV.2.s6.8.7-fr}
Soient $X$, $Y$ deux préschémas localement noethériens, $f:X\to Y$ un
morphisme localement de type fini. Alors l'ensemble des points $x\in X$ où
$f$ est lisse (ou régulier) est ouvert dans $X$.
\end{corollary}

\begin{proof}
Il résulte en effet de (0, 22.6.5) que l'ensemble des $x\in X$ vérifiant la
condition c') de (6.8.6) est ouvert dans $X$, et on conclut par (6.8.6).
\end{proof}

\begin{remark}[6.8.8]
\label{IV.2.s6.8.8-fr}
Dans (17.5.1), nous montrerons que l'équivalence de b) et c') dans (6.8.6),
ainsi que le corollaire (6.8.7), restent valables sans hypothèse noethérienne
sur $X$ et $Y$, pourvu que l'on se borne aux morphismes localement de
présentation finie.
\end{remark}

\subsection{Le théorème de platitude générique}
\label{subsection:IV.6.9}

\begin{theorem}[6.9.1]
\label{IV.2.s6.9.1-fr}
Soient $Y$ un préschéma localement noethérien et intègre, $u:X\to Y$ un
morphisme de type fini, $\mathcal F$ un $\mathcal O_X$-Module cohérent. Il
existe alors un ouvert non vide $U$ de $Y$ tel que
$\mathcal F|_{u^{-1}(U)}$ soit plat sur $U$.
\end{theorem}

\begin{proof}
On peut évidemment se borner au cas où $Y=\operatorname{Spec}(A)$ est affine;
alors $X$ est réunion finie d'ouverts affines $X_i$ de type fini sur $Y$; si,
pour chaque $i$, il y a un ouvert non vide $U_i$ dans $Y$ tel que
$\mathcal F|_{X_i\cap u^{-1}(U_i)}$ soit plat sur $U_i$, il est clair qu'en
prenant pour $U$ l'intersection des $U_i$, on répondra à la question; on peut
donc supposer que $X=\operatorname{Spec}(B)$, où $B$ est une $A$-algèbre de
type fini. On a alors $\mathcal F=\widetilde M$, où $M$ est un $B$-module de
type fini; le théorème résultera par suite du lemme suivant.
\end{proof}

\begin{lemma}[6.9.2]
\label{IV.2.s6.9.2-fr}
Soient $A$ un anneau intègre noethérien, $B$ une $A$-algèbre de type fini,
$M$ un $B$-module de type fini. Alors il existe $f\ne0$ dans $A$ tel que
$M_f$ soit un $A_f$-module libre.
\end{lemma}

\begin{proof}
